\documentclass[10pt,a4paper]{amsart}
\usepackage[margin=19mm]{geometry}
\usepackage{amsmath,amssymb,mathtools}
\usepackage[scr=rsfs,cal=euler]{mathalfa}
\usepackage{xfrac}
\usepackage[unicode,hidelinks,bookmarks=false]{hyperref}
\newcommand{\ie}{i.e.\ }
\newcommand{\cf}{cf.\ }
\newcommand{\resp}{respectively\ }
\newcommand{\Subsection}{\subsection}
\providecommand{\nobreakdash}{\nobreak\hbox{-}}
\setlength{\emergencystretch}{2em}
\multlinegap0pt
\usepackage{bm}
\usepackage{bbm}
\usepackage{dsfont}
\usepackage{xspace}
\usepackage{cancel}
\usepackage{mathtools}
\usepackage{cleveref}
\usepackage{amsthm}
\makeatletter
\@ifundefined{theorem}{\newtheorem{theorem}{Theorem}}{}
\@ifundefined{lemma}{\newtheorem{lemma}[theorem]{Lemma}}{}
\makeatother
\makeatletter
\DeclareMathOperator{\End}{End}
\DeclarePairedDelimiter{\abs}{\lvert}{\rvert}
\newcommand{\cF}{\mathcal F}
\let\emptyset\varnothing
\newcommand{\eps}{\varepsilon}
\newcommand{\ev}{\mathbb{E}}
\newcommand{\pr}{\mathbb{P}}
\expandafter\ifx\csname proof\endcsname\relax
\renewenvironment{proof}[1][\proofname]
{\par\pushQED{\qed}
	\normalfont\topsep6\p@\@plus6\p@\relax\trivlist
	\item[\hskip\labelsep\bfseries#1\@addpunct{.}]
	\ignorespaces}
{\popQED\endtrivlist\@endpefalse}
\fi
\makeatother
\title[New bounds for linear arboricity and related problems]{New bounds for linear arboricity and related problems}
\author{Micha Christoph, Nemanja Dragani\'c, Ant\'onio Gir\~ao, Eoin Hurley, Lukas Michel, Alp M\"uyesser}
\date{Source: arXiv:2507.20500v1 (CC BY 4.0)}
\begin{document}
\maketitle

\noindent\emph{Source and licence.} New bounds for linear arboricity and related problems. Version arXiv:2507.20500v1 (CC BY 4.0), distributed under the Creative Commons Attribution 4.0 International licence, as recorded in the arXiv licence metadata for this exact version and shown on the abstract page https://arxiv.org/abs/2507.20500v1. Full rights notice: licenses/arxiv-2507.20500-CC-BY-4.0.txt.\par\smallskip

\noindent\textit{Editorial scope.} Counted unit: the Lemma whose source label is \texttt{lem:lowendpointprobability} in the article (source0.tex lines 683--685, source label \texttt{lem:lowendpointprobability}), delivered together with the article's own complete original proof of it (source0.tex lines 687--707, one \texttt{proof} environment of 21 lines). No mathematics is written, repaired or re-derived: statement and proof are the article's own text line for line. Required context retained uncounted: prop:componentsize (lines 289--294). Every definition, display and label of those lines is retained unchanged. Omitted from this package and not redistributed: 1--288, 295--682, 686--686, 708--779. Nothing inside a retained excerpt is omitted or altered: every retained body is delivered exactly as the article's lines are written, so no omission receipt of any kind accompanies this package. Citations: the retained excerpts carry no citation command at all, so no bibliography entry of the article is transcribed and no reference is redistributed. Reference adaptation: none was needed and none was made; every reference inside the delivered unit resolves inside it, so no unresolved reference or citation is produced. Printed slips kept literal: no printed word, symbol, hypothesis or number of the counted statement or of its proof was altered. Layout and class: the article's own class is \texttt{article} with packages including inputenc, fontenc, mathpazo, babel, geometry, parskip, xcolor, enumitem, amsmath, amsfonts, amssymb, amsthm, mathtools, thmtools; the wrapper is the lane's pinned \texttt{amsart} editorial wrapper at 10pt on A4, which restates the theorem-like environments the retained text needs and adds the article's own macro definitions that the retained text needs (\texttt{\textbackslash End}, \texttt{\textbackslash abs}, \texttt{\textbackslash cF}, \texttt{\textbackslash emptyset}, \texttt{\textbackslash ev}, \texttt{\textbackslash pr}), with extra wrapper packages (bm, bbm, dsfont, xspace, cancel, mathtools, cleveref). The delivered numbering is the unit's own. Assets: no retained span contains a figure environment, a graphics-inclusion command, a PDF-page inclusion, a table or a TikZ picture; no image file is copied into this package, the package declares no asset dependency and no image file is redistributed with it. Licence: version arXiv:2507.20500v1 of this article is distributed under the Creative Commons Attribution 4.0 International licence, as recorded in the arXiv licence metadata for this exact version and shown on the abstract page https://arxiv.org/abs/2507.20500v1. A full rights notice is delivered at licenses/arxiv-2507.20500-CC-BY-4.0.txt. Characters: every retained excerpt is pure ASCII, so the delivered engine, XeLaTeX with its default Latin Modern text font, reports no missing character.
\begin{lemma}\label{prop:componentsize}
    Let $F$ be a minimum linear forest of a $d$-regular graph, and let $X \subseteq \End(F)$. Then,
    \[
        \abs{C(F,X)} \ge \frac{d+1}{4} \cdot \abs{\End(F) \setminus X}.
    \]
\end{lemma}

\begin{lemma}\label{lem:lowendpointprobability}
    Let $G$ be a $d$-regular graph. Then, we can choose a random minimum linear forest $F$ of $G$ in such a way that for all vertices $v \in V(G)$ it holds that $\pr(v \in \End(F)) \le 16/(d+1)$.
\end{lemma}

\begin{proof}
    We claim that for all minimum linear forests $F$ of $G$ and all $X \subseteq \End(F)$, we can choose a random linear forest $F' \in \cF(F,X)$ in such a way that for all vertices $v \in V(G)$ it holds that $\pr(v \in \End(F') \setminus X) \le 16/(d+1)$. For $X = \emptyset$, this implies the result. We prove the claim by induction on $\ell \coloneqq \abs{\End(F) \setminus X}$. The claim clearly holds for $\ell = 0$, so suppose that $\ell > 0$.

    For all $u \in C(F,X)$, let $F_u \in \cF(F,X)$ be a linear forest with $X \cup \{u\} \subseteq \End(F_u)$. By the induction hypothesis, we know that we can choose a random linear forest $F_u' \in \cF(F_u,X \cup \{u\})$ such that for all vertices $v \in V(G)$ it holds that $\pr(v \in \End(F_u') \setminus (X \cup \{u\})) \le 16/(d+1)$. Choose a random linear forest $F'' \in \cF(F,X)$ by first choosing $u \in C(F,X)$ uniformly at random and then choosing the random linear forest $F_u'$. By the linearity of expectation, we know that
    \[
        \sum_{v \in C(F,X)} \pr(v \in \End(F'') \setminus X) = \ev(\abs{\End(F'') \setminus X}) = \abs{\End(F) \setminus X}.
    \]
    By \cref{prop:componentsize}, we also have $\abs{C(F,X)} \ge \abs{\End(F) \setminus X} \cdot (d+1)/4$. This implies that the set $U \coloneqq \{v \in C(F,X) : \pr(v \in \End(F'') \setminus X) \le 8/(d+1)\}$ must contain at least half of the vertices of $C(F,X)$.

    Now, choose the random linear forest $F' \in \cF(F,X)$ by first choosing $u \in U$ uniformly at random and then choosing the random linear forest $F_u'$. For all vertices $v \in V(G) \setminus U$, we have
    \begin{align*}
        \pr(v \in \End(F') \setminus X) & = \frac{1}{\abs{U}} \sum_{u \in U} \pr(v \in \End(F_u') \setminus X) \\
        & = \frac{1}{\abs{U}} \sum_{u \in U} \pr(v \in \End(F_u') \setminus (X \cup \{u\})) \le \frac{16}{d+1}.
    \end{align*}
    For all vertices $v \in U$, we have
    \begin{align*}
        \pr(v \in \End(F') \setminus X) & = \frac{1}{\abs{U}} \sum_{u \in U} \pr(v \in \End(F_u') \setminus X) \\
        & \le \frac{1}{\abs{U}} \sum_{u \in C(F,X)} \pr(v \in \End(F_u') \setminus X) \\
        & = \frac{\abs{C(F,X)}}{\abs{U}} \cdot \pr(v \in \End(F'') \setminus X) \le \frac{16}{d+1}. \qedhere
    \end{align*}
\end{proof}

\end{document}
